\startvolume{Paper VI}{Shape composition and transfer}
\lesson{PAPER VI / 01}{Read shape through the logarithmic coordinate}{VI01_normalized_shape}
{g_\alpha(y)=\Gamma(\alpha)^{-1/2}e^{(\alpha y-e^y)/2},\qquad |g_\alpha(y)|^2\,dy=\frac{u^{\alpha-1}e^{-u}}{\Gamma(\alpha)}\,du}
{The left panel shows probability densities, the squared packet amplitudes. The right shows the same laws after $u=e^y$. Each curve has total mass one on its full domain. The logarithmic Jacobian changes the curve and its mode. Shape therefore does more than move a fixed profile along an axis.}
{For shape two, the log density is $e^{2y-e^y}$. Its derivative vanishes at $y=\log2$. The corresponding Gamma density is $u e^{-u}$, whose mode is $u=1$. The mean of the log coordinate is $1-\gamma_{\mathrm E}$, about $0.4228$, where $\gamma_{\mathrm E}$ is Euler's constant; its variance is $\pi^2/6-1$.}
{A change of coordinate preserves probability after including its Jacobian. Packet normalization uses the full real line.}
{VI Proposition 1.1, ``Normalization.'' Shapes $1,2,4,8$; finite plot windows. Omitted mass is recorded using regularized Gamma functions; no plotted curve is renormalized.}

\lesson{PAPER VI / 02}{Separate shape affinity from local curvature}{VI02_shape_affinity}
{F(\alpha,\gamma)=\frac{\Gamma((\alpha+\gamma)/2)}{\sqrt{\Gamma(\alpha)\Gamma(\gamma)}},\qquad F(\alpha,\alpha+d)=1-\frac{\psi_1(\alpha)}8d^2+O(d^3)}
{The first panel compares every displayed second shape with the fixed shape two. Only the equal-shape point reaches one. The second enlarges a neighborhood of that point and compares the exact affinity with its quadratic expansion. Their separation away from zero shows why a local expansion should not be used as a global identity.}
{For shapes two and four, $F=2/\sqrt6=\sqrt{2/3}$. Thus the squared overlap at zero translation and modulation is exactly $2/3$. At increment $d=0.2$, the exact affinity is about $0.996969$; the quadratic approximation gives $0.996775$. Its curvature comes from $\psi_1(2)=\pi^2/6-1$.}
{Origin affinity compares shapes at one common coordinate. A later translation can increase their overlap.}
{VI Proposition 2.2, ``Strict shape affinity,'' and its following local expansion. The local quadratic is displayed as an approximation; the shape metric coefficient is $\psi_1(\alpha)/4$.}

\lesson{PAPER VI / 03}{Keep magnitude and phase in cross ambiguity}{VI03_cross_ambiguity}
{\mathcal A_{\alpha,\gamma}(h,\eta)=e^{(\alpha-\gamma)h/4}\frac{\Gamma(\lambda+i\eta)}{\sqrt{\Gamma(\alpha)\Gamma(\gamma)}}\operatorname{sech}^{\lambda+i\eta}(h/2),\quad\lambda=(\alpha+\gamma)/2}
{The two panels represent one complex overlap. Magnitude measures localization; the principal phase records information that a squared magnitude removes. A phase-color jump is a branch wrap, not a zero. Unequal shapes shift the magnitude peak away from zero displacement, while the frequency envelope remains symmetric. The white cross marks the zero-frequency displacement maximum.}
{Take $\alpha=2$, $\gamma=4$. At zero frequency, the magnitude is $\sqrt{2/3}\,e^{-h/2}\operatorname{sech}^3(h/2)$. Differentiating its logarithm gives the maximum at $h=-\log2$, with value about $0.967700$. The origin value is only $\sqrt{2/3}$. At $(h,\eta)=(0,1)$ the squared magnitude is $5\pi/(3\sinh\pi)$.}
{Translation, modulation, and shape are distinct coordinates. The complex Gamma factor retains their phase relation.}
{VI Theorem 2.1, ``Cross-shape Gamma--sech formula.'' Shapes $2,4$ on $[-3,3]^2$; exact-formula samples. Independent overlap quadrature is checked separately.}

\lesson{PAPER VI / 04}{Factor energy into two independent laws}{VI04_factorization}
{\frac{|\mathcal A_{\alpha,\gamma}(h,\eta)|^2}{2\pi}=p_{\alpha,\gamma}(h)\rho_{\alpha+\gamma}(\eta),\qquad \rho_s(\eta)=\frac{2^{s-1}|\Gamma(s/2+i\eta)|^2}{\pi\Gamma(s)}}
{The displacement panel reflects when the two shapes are exchanged. The frequency panel uses only their sum. These are probability marginals of the normalized cross-spectrogram, whose product is its full energy density. Independence here belongs to that product probability law. It does not say that the complex ambiguity factors into two real densities or loses its phase.}
{With shapes two and four, displacement is the logit of a Beta$(2,4)$ variable. Its mean is $\psi(2)-\psi(4)=-5/6$, while the symmetric shape-six frequency has mean zero and variance $3/2$. At the origin, $p_{2,4}(0)=5/16$ and $\rho_6(0)=16/(15\pi)$, so the joint density is exactly $1/(3\pi)$.}
{A positive two-dimensional energy law can have independent marginals even when its underlying amplitude has nontrivial phase.}
{VI Theorem 3.1, ``Exact cross-spectrogram factorization.'' Position tails use Beta CDFs; frequency tails use a stated Meixner Chernoff bound. Finite-window mass is not reset to one.}

\lesson{PAPER VI / 05}{Add shapes by convolution}{VI05_shape_convolution}
{\rho_{s+t}=\rho_s*\rho_t,\qquad \rho_{s+2}(x)=\frac{s^2+4x^2}{s(s+1)}\rho_s(x)}
{The first panel overlays a closed density with an independently integrated convolution. The second uses a different operation: pointwise quadratic multiplication raises shape by two. Both identities change shape, but their inputs and computational costs differ. The convolution adds independent random variables. The ladder changes one density locally at each displayed frequency.}
{Shape one has density $\operatorname{sech}(\pi x)$. Convolving two such laws gives shape two, $2x/\sinh(\pi x)$, with the removable origin value $2/\pi$. Its variance is $1/4+1/4=1/2$. Separately, multiplying shape one by $(1+4x^2)/2$ produces shape three. At $x=1$, shape two has density about $0.173179$.}
{The characteristic function $\operatorname{sech}^s(h/2)$ turns shape addition into multiplication. The Gamma recurrence gives the ladder.}
{VI Theorem 4.1, ``Meixner shape semigroup,'' and Proposition 4.2, ``Two-step density ladder.'' Convolution integrates on $[-8,8]$; its omitted integral is uniformly bounded by $1.55\times10^{-11}$.}

\lesson{PAPER VI / 06}{Continuous shape produces jumps and cumulants}{VI06_levy_bernoulli}
{-\log\cosh(\xi/2)=\int_0^\infty\frac{\cos(\xi x)-1}{x\sinh(\pi x)}\,dx,\qquad \kappa_{2m}=s\frac{(2^{2m}-1)|B_{2m}|}{2m}}
{The generator panel compares its exact Fourier symbol with the L\'evy integral. The numerator cancels the apparent origin singularity, so the combined integrand is finite there. The second panel tracks even cumulants through shape. The second and sixth cumulants overlap because both equal $s/4$; their coincidence does not make the distribution Gaussian.}
{At shape four, variance is one, the fourth cumulant is $1/2$, and the sixth is one. A Gaussian with variance one would have both higher cumulants zero. The nonzero values measure departures from Gaussianity. For dual frequency $\xi=2$, the unit-shape generator symbol is $-\log\cosh1$, about $-0.433781$.}
{The L\'evy measure describes infinitesimal shape increments; Bernoulli numbers encode the complete even-cumulant sequence.}
{VI Theorems 5.1 and 5.2, ``L\'evy--Khintchine representation'' and ``Bernoulli cumulant law.'' Quadrature uses $[0,10]$, stable sine-squared cancellation, and a tail bound below $2.90\times10^{-15}$.}

\lesson{PAPER VI / 07}{Split one polynomial without losing energy}{VI07_two_child_refinement}
{p_{n,\alpha+\gamma}(x+y)=\sum_{r=0}^n C_{n,r}^{\alpha,\gamma}p_{r,\alpha}(x)p_{n-r,\gamma}(y),\qquad (C_{n,r}^{\alpha,\gamma})^2=\binom nr\frac{(\alpha)_r(\gamma)_{n-r}}{(\alpha+\gamma)_n}}
{Every bar is an allowed degree allocation between two children. Their heights sum to one, so the transfer preserves coefficient energy. The polynomial panel overlays the parent with the sum of all five products. Individual products can be large or signed; the probability bars describe squared transfer coefficients, not pointwise polynomial contributions.}
{Use degree four and child shapes two and three. The energy weights, in order $r=0,\ldots,4$, are $(15,20,18,12,5)/70$. Their sum is one and their mean allocation to the first child is $8/5$. At $x=0.6$, $y=0.85$, the normalized parent is approximately $-0.716400$, reproduced by the complete five-term child sum.}
{An isometric refinement is a complete change of polynomial coordinates. Omitting an allocation changes the reconstruction.}
{VI Theorem 6.1, ``Two-child packet refinement.'' Degree four, shapes $2,3$, fixed $y=0.85$; all five terms are retained. Numerical pointwise reconstruction and coefficient normalization are checked.}

\lesson{PAPER VI / 08}{Choose a refinement tree freely}{VI08_coassociative_refinement}
{C_{n,\mathbf r}^{\boldsymbol\alpha}=\sqrt{\frac{n!\prod_j(\alpha_j)_{r_j}}{(\sum_j\alpha_j)_n\prod_jr_j!}},\qquad C_{n,r_1+r_2}^{\alpha_1+\alpha_2,\alpha_3}C_{r_1+r_2,r_1}^{\alpha_1,\alpha_2}=C_{n,\mathbf r}^{\boldsymbol\alpha}}
{The triangle contains every degree-four allocation among three children. Color encodes energy; marker area increases with it. The second panel lists all fifteen terminal amplitudes and overlays two binary refinement orders. Coincident markers express coassociativity. The geometric arrangement is a degree-allocation simplex, not the phase space used by the ambiguity function.}
{Take shapes $(2,3,1)$ and allocation $(1,2,1)$. The direct squared coefficient is $2/21$. Splitting the first two children first gives squared factors $5/18$ and $12/35$. Splitting the last two first gives $20/63$ and $3/10$. Both products equal $2/21$, and both amplitudes are the positive square root, about $0.308607$.}
{Intermediate shape factors cancel exactly. A deeper binary tree gives the same final many-child coefficients.}
{VI Theorem 7.1, ``Multinomial packet refinement,'' and Corollary 7.2, ``Coassociativity.'' All fifteen compositions of degree four are shown; both trees and total energy are checked.}

\lesson{PAPER VI / 09}{Read conditional transfer from its singular values}{VI09_martingale_spectrum}
{T_{s,t}p_{n,s+t}=\sigma_n(s,t)p_{n,s},\qquad \sigma_n(s,t)=\sqrt{\frac{(s)_n}{(s+t)_n}},\quad \sigma_n(s,t)\sim\sqrt{\frac{\Gamma(s+t)}{\Gamma(s)}}n^{-t/2}}
{Each spectral curve measures contraction of one polynomial degree. Larger increments suppress high degrees more strongly. Degree zero, excluded from the logarithmic horizontal axis, remains exactly one. The polynomial panel shows why orthonormal packets contract even though the standard, unnormalized Meixner--Pollaczek polynomials satisfy the martingale identity without a scalar loss.}
{For $s=2$, $t=3$, the first singular values are $1$, $\sqrt{2/5}$, $1/\sqrt5$, and $2/\sqrt{35}$. Since a shape-three increment has variance $3/4$, conditioning $(2(X+Y)^2-5/2)/\sqrt{15}$ gives $(2X^2-1)/\sqrt{15}$. This equals $p_{2,2}(X)/\sqrt5$, and at $X=1$ it is $1/\sqrt{15}$.}
{The exact cocycle permits successive increments. Compactness holds for every positive increment; Schatten membership requires $pt>2$.}
{VI Theorem 8.1, ``Polynomial martingale law,'' and Proposition 8.2, ``Cocycle and compactness.'' Degrees through 60; a rank-61 truncation at increment three has exact omitted norm $\sigma_{61}\approx0.00957$.}

\lesson{PAPER VI / 10}{A continuous coherent orbit resolves identity}{VI10_coherent_frame}
{\psi_{\alpha,u,\eta}(y)=e^{i\eta(y-u/2)}g_\alpha(y-u),\qquad \frac1{2\pi}\int_{\mathbb R^2}|\langle\psi_{\alpha,u,\eta},f\rangle|^2\,du\,d\eta=\|f\|_2^2}
{Translation moves the window; modulation oscillates its complex amplitude. Only real parts are drawn on the left. The right integrates the coherent-transform energy of $f=g_2$ over a position strip, including every modulation frequency. Its approach to one illustrates the continuous frame normalization. Three displayed states alone do not form the identity resolution.}
{For $f=g_2$, the frame energy density is $p_{2,2}(u)\rho_4(\eta)$. Integrating frequency removes its normalized factor. In the strip $|u|\le\log3$, the logistic endpoints are $1/4$ and $3/4$. The Beta$(2,2)$ CDF is $3v^2-2v^3$, so the captured energy is exactly $11/16$.}
{An overlap kernel becomes a coherent transform and a reproducing kernel when the full normalized orbit is used.}
{VI Theorems 9.1 and 9.2, ``Cross-shape coherent overlap'' and ``Continuous tight frame.'' Shape two; strip energy uses an exact Beta CDF with frequency integrated on all $\mathbb R$.}

\lesson{PAPER VI / 11}{Positive smoothing can have unstable inversion}{VI11_berezin_inverse}
{\widehat{\mathcal B}_\alpha(\xi,x)=\frac{|\Gamma(\alpha+i\xi)|^2}{\Gamma(\alpha)^2}\operatorname{sech}^{2\alpha}(x/2)>0,\qquad \widetilde F=\mathcal B_\alpha*F}
{The multiplier stays positive while shrinking quickly along either dual coordinate. The bars show its action on four cosine modes: convolution preserves each mode and multiplies its amplitude by the displayed number. Therefore exact information survives algebraically, but recovering a small attenuated mode requires amplifying any error attached to it.}
{At shape two, $\widehat{\mathcal B}_2(1,0)=2\pi/\sinh\pi\approx0.544058$. The mode $(3,0)$ retains only $30\pi/\sinh(3\pi)\approx0.015212$, requiring inverse amplification about $65.7$. For mode $(1,2)$, the additional factor is $\operatorname{sech}^4(1)$, reducing its amplitude to about $0.095960$. The constant mode retains amplitude one.}
{Injectivity is an exact uniqueness statement. Stable numerical recovery requires an additional regularization or restricted function class.}
{VI Theorem 10.1, ``Exact Berezin factorization,'' and Corollary 10.2, ``Injectivity and ill-posed inversion.'' Bounded cosine symbols illustrate attenuation; the $L^2$ inverse statement is the source theorem.}

\lesson{PAPER VI / 12}{Rescale shape toward a Gaussian heat kernel}{VI12_gaussian_heat}
{Z_\alpha=\sqrt{\alpha/2}\,H_\alpha,\quad W_\alpha=\sqrt{2/\alpha}\,\Omega_\alpha\ \Longrightarrow\ N(0,1),\qquad S_\alpha F\longrightarrow e^{\Delta/2}F}
{The two panels use different rescalings because unscaled displacement narrows while frequency broadens. Both normalized laws approach the same standard normal curve. At finite shape, the position variance is still slightly above one; frequency variance is exactly one. Matching a limiting curve does not remove the higher finite-shape cumulants or establish a function-space convergence theorem.}
{At shape two, $\operatorname{Var}(Z_2)=2\psi_1(2)\approx1.289868$, while $\operatorname{Var}(W_2)=1$. At shape 32, the first variance is about $1.015788$. The rescaled two-dimensional multiplier at dual mode $(1,1)$ is then about $0.365956$, approaching the heat multiplier $e^{-1}$. For a Gaussian test symbol $e^{-|z|^2/2}$, the limiting smoothed value is $\tfrac12e^{-|z|^2/4}$.}
{The natural rescaling gives independent standard-normal marginals and the two-dimensional heat operator as the large-shape limit.}
{VI Theorem 11.1, ``Gaussian marginal limits,'' and Corollary 11.2, ``Heat-operator limit.'' Exact finite-shape densities; tail bounds recorded. All-Schwartz-seminorm convergence belongs to the source proof, not the plotted comparison.}
